\chapter{Theoretical Background}

\vspace*{0.3cm}

This chapter presents the main theoretical concepts used in the project. The work
is based on three connected ideas: Small Language Models, fine-tuning, and
optimization for Edge AI constraints. Understanding these concepts is necessary
before describing the implementation and the experimental comparison.

\section{Small Language Models}

Small Language Models are language models designed with a reduced number of
parameters compared with Large Language Models. Their objective is to provide
useful natural language processing capabilities while requiring fewer
computational resources. This makes them more suitable for limited environments,
such as personal computers, embedded systems, mobile devices, and edge platforms.

\vspace*{0.15cm}

In this project, the selected model is \texttt{google/flan-t5-small}. It belongs
to the T5 family, where many language tasks are represented as text-to-text
problems. The FLAN version is instruction-tuned, which means that it has already
been trained to follow natural language instructions. This makes it a good
candidate for experiments based on instruction-following data.

\section{Fine-Tuning}

Fine-tuning is the process of adapting a pre-trained model to a specific task or
dataset. Instead of training a model from the beginning, fine-tuning starts from
existing learned parameters and updates them using a smaller target dataset. This
approach reduces training cost and usually improves performance on the target
task.

\vspace*{0.15cm}

For instruction tuning, each training example contains an input prompt and an
expected answer. In this project, the prompt is stored in the \texttt{source\_text}
column, while the expected answer is stored in the \texttt{target\_text} column.
The model learns to generate the target answer from the given instruction-based
source text.

\section{Edge AI Constraints}

Edge AI consists of deploying artificial intelligence models directly on local
devices instead of relying only on cloud computing. This can improve latency,
privacy, autonomy, and availability. However, edge devices generally have limited
memory, processing power, and energy capacity.

\vspace*{0.15cm}

Because of these constraints, the evaluation of a language model should not be
limited to accuracy or loss. It is also important to consider model size,
inference time, latency, memory consumption, and energy usage. These criteria help
determine whether a model is practical for real deployment on constrained
hardware.

\section{Optimization in Deep Learning}

Optimization algorithms are used to update the model parameters during training.
The objective is to minimize the loss function, which measures the difference
between the model prediction and the expected output. The choice of optimizer can
strongly affect training speed, stability, memory usage, and final performance.

\vspace*{0.15cm}

This project studies several optimization methods. Stochastic Gradient Descent is
a simple and classical method based on gradient updates. SGD with momentum
improves the basic method by accumulating part of the previous update direction.
Adam uses adaptive learning rates and is widely used for transformer training.
L-BFGS is a quasi-Newton method that approximates second-order information, but it
is usually expensive for large neural networks. Newton-CG is discussed
theoretically because it requires second-order information that is not practical
for this fine-tuning pipeline.

\section{Model Compression Perspective}

Model compression aims to reduce the size and computational cost of a model while
preserving as much performance as possible. Common compression techniques include
quantization, pruning, and knowledge distillation. Although this project mainly
focuses on fine-tuning and optimizer comparison, the collected metrics provide a
baseline for future compression experiments.

\vspace*{0.15cm}

The theoretical foundation of this work is therefore centered on finding a balance
between model quality and efficiency. This balance is essential when preparing
Small Language Models for Edge AI applications.
